% preamble.tex
% Shared package imports and environment definitions for the course notes.
% This file is intended to be loaded from main.tex via
%     % preamble.tex
% Shared package imports and environment definitions for the course notes.
% This file is intended to be loaded from main.tex via
%     % preamble.tex
% Shared package imports and environment definitions for the course notes.
% This file is intended to be loaded from main.tex via
%     % preamble.tex
% Shared package imports and environment definitions for the course notes.
% This file is intended to be loaded from main.tex via
%     \input{preamble}

% --------------------------------------------------------------------------
% Encoding, fonts, and typography
% --------------------------------------------------------------------------

\usepackage[T1]{fontenc}
\usepackage[utf8]{inputenc} % Harmless on modern LaTeX; useful on older systems.
\usepackage{lmodern}
\usepackage{microtype}

% --------------------------------------------------------------------------
% Page layout and document structure
% --------------------------------------------------------------------------

\usepackage[
  letterpaper,
  inner=1.25in,
  outer=1.00in,
  top=1.00in,
  bottom=1.10in,
  headheight=15pt
]{geometry}

\usepackage{setspace}
\setstretch{1.05}

\usepackage{emptypage}
\usepackage{fancyhdr}

\pagestyle{fancy}
\fancyhf{}
\fancyhead[LE]{\small\thepage}
\fancyhead[RO]{\small\thepage}
\fancyhead[LO]{\small\nouppercase{\rightmark}}
\fancyhead[RE]{\small\nouppercase{\leftmark}}
\renewcommand{\headrulewidth}{0.4pt}
\renewcommand{\footrulewidth}{0pt}

\fancypagestyle{plain}{
  \fancyhf{}
  \fancyfoot[C]{\small\thepage}
  \renewcommand{\headrulewidth}{0pt}
}

\usepackage{titlesec}

\titleformat{\chapter}[display]
  {\normalfont\bfseries}
  {\Large Chapter \thechapter}
  {0.5ex}
  {\Huge}

\titleformat{name=\chapter,numberless}[display]
  {\normalfont\bfseries}
  {}
  {0pt}
  {\Huge}

\titlespacing*{\chapter}{0pt}{-10pt}{30pt}

% --------------------------------------------------------------------------
% Mathematics
% --------------------------------------------------------------------------

\usepackage{amsmath}
\usepackage{amssymb}
\usepackage{amsthm}
\usepackage{mathtools}
\usepackage{mathrsfs}
\usepackage{bm}
\usepackage{dsfont}

\numberwithin{equation}{chapter}

% --------------------------------------------------------------------------
% Lists, tables, and general utilities
% --------------------------------------------------------------------------

\usepackage{enumitem}
\setlist{
  itemsep=2pt,
  topsep=4pt,
  parsep=0pt
}

\usepackage{booktabs}
\usepackage{array}
\usepackage{longtable}
\usepackage{tabularx}
\usepackage{multicol}
\usepackage{ragged2e}

\usepackage{xcolor}
\usepackage{xparse}
\usepackage{etoolbox}
\usepackage{ifthen}
\usepackage{comment}

% --------------------------------------------------------------------------
% Graphics, diagrams, and commutative diagrams
% --------------------------------------------------------------------------

\usepackage{graphicx}
\graphicspath{{figures/}}

\usepackage{tikz}
\usetikzlibrary{
  arrows.meta,
  calc,
  decorations.markings,
  intersections,
  matrix,
  patterns,
  positioning,
  shapes.geometric
}

\usepackage{tikz-cd}

% Basic TikZ styles for fans, lattices, and polytopes.
\tikzset{
  lattice point/.style={
    circle,
    fill,
    inner sep=1.4pt
  },
  ray/.style={
    -{Latex[length=2.2mm]},
    thick
  },
  fan cone/.style={
    fill=black!8,
    draw=black!55
  },
  polytope/.style={
    thick
  },
  toric divisor/.style={
    thick,
    dashed
  }
}

% --------------------------------------------------------------------------
% Color palette
% --------------------------------------------------------------------------

\definecolor{DictionaryBlue}{HTML}{1F4E79}
\definecolor{SurfaceGreen}{HTML}{3F6B4F}
\definecolor{WarningRed}{HTML}{8B2E2E}
\definecolor{ExampleGold}{HTML}{8A641B}
\definecolor{LightGray}{HTML}{F4F4F4}
\definecolor{PastelYellow}{HTML}{FFF8CC}
\definecolor{PausePurple}{HTML}{6F5A8A}
\definecolor{PredictionTeal}{HTML}{287D7A}
\definecolor{CheatsheetSlate}{HTML}{536878}

% --------------------------------------------------------------------------
% Boxed pedagogical environments
% --------------------------------------------------------------------------

\usepackage[most]{tcolorbox}

\tcbset{
  enhanced,
  breakable,
  boxrule=0.7pt,
  arc=1.5mm,
  left=2mm,
  right=2mm,
  top=1.5mm,
  bottom=1.5mm,
  before skip=8pt,
  after skip=8pt
}

\newtcolorbox{chapteroverview}{
  colback=PastelYellow,
  colframe=ExampleGold!70!black,
  title={Chapter overview},
  fonttitle=\bfseries
}

\newtcolorbox{threedictionaries}{
  colback=DictionaryBlue!4,
  colframe=DictionaryBlue,
  title={The three dictionaries},
  fonttitle=\bfseries
}

\newtcolorbox{surfacelaboratory}{
  colback=SurfaceGreen!5,
  colframe=SurfaceGreen,
  title={Surface laboratory},
  fonttitle=\bfseries
}

\newtcolorbox{warningbox}{
  colback=WarningRed!4,
  colframe=WarningRed,
  title={Warning},
  fonttitle=\bfseries
}

\newtcolorbox{guidingquestion}{
  colback=ExampleGold!5,
  colframe=ExampleGold,
  title={Guiding question},
  fonttitle=\bfseries
}

\newtcolorbox{pauseandthink}{
  colback=PausePurple!5,
  colframe=PausePurple,
  title={Pause and think},
  fonttitle=\bfseries
}

\newtcolorbox{prediction}{
  colback=PredictionTeal!5,
  colframe=PredictionTeal,
  title={Prediction},
  fonttitle=\bfseries
}

\newtcolorbox{dictionarytable}{
  colback=DictionaryBlue!3,
  colframe=DictionaryBlue,
  title={Dictionary summary},
  fonttitle=\bfseries
}

\newtcolorbox{agcheatsheet}{
  colback=CheatsheetSlate!5,
  colframe=CheatsheetSlate,
  title={Algebraic geometry cheatsheet},
  fonttitle=\bfseries
}

% --------------------------------------------------------------------------
% Theorem environments
% --------------------------------------------------------------------------

\theoremstyle{plain}
\newtheorem{theorem}{Theorem}[chapter]
\newtheorem{proposition}[theorem]{Proposition}
\newtheorem{lemma}[theorem]{Lemma}
\newtheorem{corollary}[theorem]{Corollary}

\theoremstyle{definition}
\newtheorem{definition}[theorem]{Definition}
\newtheorem{construction}[theorem]{Construction}
\newtheorem{example}[theorem]{Example}
\newtheorem{exercise}[theorem]{Exercise}
\newtheorem{question}[theorem]{Question}

\theoremstyle{remark}
\newtheorem{remark}[theorem]{Remark}
\newtheorem{observation}[theorem]{Observation}
\newtheorem{convention}[theorem]{Convention}
\newtheorem{notation}[theorem]{Notation}

% Unnumbered variants.
\newtheorem*{theorem*}{Theorem}
\newtheorem*{proposition*}{Proposition}
\newtheorem*{definition*}{Definition}
\newtheorem*{example*}{Example}
\newtheorem*{remark*}{Remark}

% --------------------------------------------------------------------------
% Exercises and solutions
% --------------------------------------------------------------------------

\newcounter{exerciseblock}[chapter]
\renewcommand{\theexerciseblock}{\thechapter.\arabic{exerciseblock}}

\NewDocumentEnvironment{exercises}{O{Exercises}}
{
  \section*{#1}
  \addcontentsline{toc}{section}{#1}
  \begin{enumerate}[label=\textbf{\thechapter.\arabic*.},wide=0pt]
}
{
  \end{enumerate}
}

\newif\ifshowsolutions
\showsolutionsfalse
% To compile instructor solutions, change the previous line to:
% \showsolutionstrue

\NewDocumentEnvironment{solution}{}
{
  \ifshowsolutions
    \par\smallskip
    \noindent\textbf{Solution. }\ignorespaces
  \else
    \expandafter\comment
  \fi
}
{
  \ifshowsolutions
    \par\smallskip
  \else
    \expandafter\endcomment
  \fi
}

% --------------------------------------------------------------------------
% Cross-references
% --------------------------------------------------------------------------

\usepackage{hyperref}

\hypersetup{
  colorlinks=true,
  linkcolor=DictionaryBlue,
  citecolor=SurfaceGreen,
  urlcolor=WarningRed,
  pdfauthor={},
  pdftitle={The Three Dictionaries of Toric Geometry},
  pdfsubject={Notes for Renzo's Fall 2026 course at Colorado State University},
  bookmarksopen=true,
  bookmarksnumbered=true
}

\usepackage[nameinlink,noabbrev,capitalise]{cleveref}

\crefname{theorem}{theorem}{theorems}
\crefname{proposition}{proposition}{propositions}
\crefname{lemma}{lemma}{lemmas}
\crefname{corollary}{corollary}{corollaries}
\crefname{definition}{definition}{definitions}
\crefname{construction}{construction}{constructions}
\crefname{example}{example}{examples}
\crefname{exercise}{exercise}{exercises}
\crefname{question}{question}{questions}
\crefname{remark}{remark}{remarks}
\crefname{observation}{observation}{observations}
\crefname{convention}{convention}{conventions}
\crefname{notation}{notation}{notations}

% --------------------------------------------------------------------------
% Bibliography
% --------------------------------------------------------------------------

\usepackage[
  backend=biber,
  style=alphabetic,
  sorting=nyt,
  maxbibnames=99
]{biblatex}

% The present notes use chapter-level prose references rather than citation
% commands.  This guard keeps the project ready for a future bibliography
% without requiring an empty references.bib file in the current release.
\IfFileExists{references.bib}{\addbibresource{references.bib}}{}

% --------------------------------------------------------------------------
% Index
% --------------------------------------------------------------------------

\usepackage{makeidx}
\makeindex

% --------------------------------------------------------------------------
% Drafting tools
% --------------------------------------------------------------------------

% Margin notes can be used during development:
%     \marginnote{Check this proof.}
\usepackage{marginnote}

\newcommand{\todo}[1]{%
  \marginpar{\raggedright\scriptsize\color{WarningRed}\textbf{TODO:} #1}%
}

% Set to true while drafting to display labels in the margins.
\newif\ifshowlabels
\showlabelsfalse

\ifshowlabels
  \usepackage{showkeys}
\fi

% --------------------------------------------------------------------------
% Miscellaneous formatting
% --------------------------------------------------------------------------

\allowdisplaybreaks
\setcounter{tocdepth}{2}
\setcounter{secnumdepth}{3}

\renewcommand{\qedsymbol}{\(\square\)}

% Prevent isolated lines at page boundaries.
\widowpenalty=10000
\clubpenalty=10000
\displaywidowpenalty=10000


% --------------------------------------------------------------------------
% Encoding, fonts, and typography
% --------------------------------------------------------------------------

\usepackage[T1]{fontenc}
\usepackage[utf8]{inputenc} % Harmless on modern LaTeX; useful on older systems.
\usepackage{lmodern}
\usepackage{microtype}

% --------------------------------------------------------------------------
% Page layout and document structure
% --------------------------------------------------------------------------

\usepackage[
  letterpaper,
  inner=1.25in,
  outer=1.00in,
  top=1.00in,
  bottom=1.10in,
  headheight=15pt
]{geometry}

\usepackage{setspace}
\setstretch{1.05}

\usepackage{emptypage}
\usepackage{fancyhdr}

\pagestyle{fancy}
\fancyhf{}
\fancyhead[LE]{\small\thepage}
\fancyhead[RO]{\small\thepage}
\fancyhead[LO]{\small\nouppercase{\rightmark}}
\fancyhead[RE]{\small\nouppercase{\leftmark}}
\renewcommand{\headrulewidth}{0.4pt}
\renewcommand{\footrulewidth}{0pt}

\fancypagestyle{plain}{
  \fancyhf{}
  \fancyfoot[C]{\small\thepage}
  \renewcommand{\headrulewidth}{0pt}
}

\usepackage{titlesec}

\titleformat{\chapter}[display]
  {\normalfont\bfseries}
  {\Large Chapter \thechapter}
  {0.5ex}
  {\Huge}

\titleformat{name=\chapter,numberless}[display]
  {\normalfont\bfseries}
  {}
  {0pt}
  {\Huge}

\titlespacing*{\chapter}{0pt}{-10pt}{30pt}

% --------------------------------------------------------------------------
% Mathematics
% --------------------------------------------------------------------------

\usepackage{amsmath}
\usepackage{amssymb}
\usepackage{amsthm}
\usepackage{mathtools}
\usepackage{mathrsfs}
\usepackage{bm}
\usepackage{dsfont}

\numberwithin{equation}{chapter}

% --------------------------------------------------------------------------
% Lists, tables, and general utilities
% --------------------------------------------------------------------------

\usepackage{enumitem}
\setlist{
  itemsep=2pt,
  topsep=4pt,
  parsep=0pt
}

\usepackage{booktabs}
\usepackage{array}
\usepackage{longtable}
\usepackage{tabularx}
\usepackage{multicol}
\usepackage{ragged2e}

\usepackage{xcolor}
\usepackage{xparse}
\usepackage{etoolbox}
\usepackage{ifthen}
\usepackage{comment}

% --------------------------------------------------------------------------
% Graphics, diagrams, and commutative diagrams
% --------------------------------------------------------------------------

\usepackage{graphicx}
\graphicspath{{figures/}}

\usepackage{tikz}
\usetikzlibrary{
  arrows.meta,
  calc,
  decorations.markings,
  intersections,
  matrix,
  patterns,
  positioning,
  shapes.geometric
}

\usepackage{tikz-cd}

% Basic TikZ styles for fans, lattices, and polytopes.
\tikzset{
  lattice point/.style={
    circle,
    fill,
    inner sep=1.4pt
  },
  ray/.style={
    -{Latex[length=2.2mm]},
    thick
  },
  fan cone/.style={
    fill=black!8,
    draw=black!55
  },
  polytope/.style={
    thick
  },
  toric divisor/.style={
    thick,
    dashed
  }
}

% --------------------------------------------------------------------------
% Color palette
% --------------------------------------------------------------------------

\definecolor{DictionaryBlue}{HTML}{1F4E79}
\definecolor{SurfaceGreen}{HTML}{3F6B4F}
\definecolor{WarningRed}{HTML}{8B2E2E}
\definecolor{ExampleGold}{HTML}{8A641B}
\definecolor{LightGray}{HTML}{F4F4F4}
\definecolor{PastelYellow}{HTML}{FFF8CC}
\definecolor{PausePurple}{HTML}{6F5A8A}
\definecolor{PredictionTeal}{HTML}{287D7A}
\definecolor{CheatsheetSlate}{HTML}{536878}

% --------------------------------------------------------------------------
% Boxed pedagogical environments
% --------------------------------------------------------------------------

\usepackage[most]{tcolorbox}

\tcbset{
  enhanced,
  breakable,
  boxrule=0.7pt,
  arc=1.5mm,
  left=2mm,
  right=2mm,
  top=1.5mm,
  bottom=1.5mm,
  before skip=8pt,
  after skip=8pt
}

\newtcolorbox{chapteroverview}{
  colback=PastelYellow,
  colframe=ExampleGold!70!black,
  title={Chapter overview},
  fonttitle=\bfseries
}

\newtcolorbox{threedictionaries}{
  colback=DictionaryBlue!4,
  colframe=DictionaryBlue,
  title={The three dictionaries},
  fonttitle=\bfseries
}

\newtcolorbox{surfacelaboratory}{
  colback=SurfaceGreen!5,
  colframe=SurfaceGreen,
  title={Surface laboratory},
  fonttitle=\bfseries
}

\newtcolorbox{warningbox}{
  colback=WarningRed!4,
  colframe=WarningRed,
  title={Warning},
  fonttitle=\bfseries
}

\newtcolorbox{guidingquestion}{
  colback=ExampleGold!5,
  colframe=ExampleGold,
  title={Guiding question},
  fonttitle=\bfseries
}

\newtcolorbox{pauseandthink}{
  colback=PausePurple!5,
  colframe=PausePurple,
  title={Pause and think},
  fonttitle=\bfseries
}

\newtcolorbox{prediction}{
  colback=PredictionTeal!5,
  colframe=PredictionTeal,
  title={Prediction},
  fonttitle=\bfseries
}

\newtcolorbox{dictionarytable}{
  colback=DictionaryBlue!3,
  colframe=DictionaryBlue,
  title={Dictionary summary},
  fonttitle=\bfseries
}

\newtcolorbox{agcheatsheet}{
  colback=CheatsheetSlate!5,
  colframe=CheatsheetSlate,
  title={Algebraic geometry cheatsheet},
  fonttitle=\bfseries
}

% --------------------------------------------------------------------------
% Theorem environments
% --------------------------------------------------------------------------

\theoremstyle{plain}
\newtheorem{theorem}{Theorem}[chapter]
\newtheorem{proposition}[theorem]{Proposition}
\newtheorem{lemma}[theorem]{Lemma}
\newtheorem{corollary}[theorem]{Corollary}

\theoremstyle{definition}
\newtheorem{definition}[theorem]{Definition}
\newtheorem{construction}[theorem]{Construction}
\newtheorem{example}[theorem]{Example}
\newtheorem{exercise}[theorem]{Exercise}
\newtheorem{question}[theorem]{Question}

\theoremstyle{remark}
\newtheorem{remark}[theorem]{Remark}
\newtheorem{observation}[theorem]{Observation}
\newtheorem{convention}[theorem]{Convention}
\newtheorem{notation}[theorem]{Notation}

% Unnumbered variants.
\newtheorem*{theorem*}{Theorem}
\newtheorem*{proposition*}{Proposition}
\newtheorem*{definition*}{Definition}
\newtheorem*{example*}{Example}
\newtheorem*{remark*}{Remark}

% --------------------------------------------------------------------------
% Exercises and solutions
% --------------------------------------------------------------------------

\newcounter{exerciseblock}[chapter]
\renewcommand{\theexerciseblock}{\thechapter.\arabic{exerciseblock}}

\NewDocumentEnvironment{exercises}{O{Exercises}}
{
  \section*{#1}
  \addcontentsline{toc}{section}{#1}
  \begin{enumerate}[label=\textbf{\thechapter.\arabic*.},wide=0pt]
}
{
  \end{enumerate}
}

\newif\ifshowsolutions
\showsolutionsfalse
% To compile instructor solutions, change the previous line to:
% \showsolutionstrue

\NewDocumentEnvironment{solution}{}
{
  \ifshowsolutions
    \par\smallskip
    \noindent\textbf{Solution. }\ignorespaces
  \else
    \expandafter\comment
  \fi
}
{
  \ifshowsolutions
    \par\smallskip
  \else
    \expandafter\endcomment
  \fi
}

% --------------------------------------------------------------------------
% Cross-references
% --------------------------------------------------------------------------

\usepackage{hyperref}

\hypersetup{
  colorlinks=true,
  linkcolor=DictionaryBlue,
  citecolor=SurfaceGreen,
  urlcolor=WarningRed,
  pdfauthor={},
  pdftitle={The Three Dictionaries of Toric Geometry},
  pdfsubject={Notes for Renzo's Fall 2026 course at Colorado State University},
  bookmarksopen=true,
  bookmarksnumbered=true
}

\usepackage[nameinlink,noabbrev,capitalise]{cleveref}

\crefname{theorem}{theorem}{theorems}
\crefname{proposition}{proposition}{propositions}
\crefname{lemma}{lemma}{lemmas}
\crefname{corollary}{corollary}{corollaries}
\crefname{definition}{definition}{definitions}
\crefname{construction}{construction}{constructions}
\crefname{example}{example}{examples}
\crefname{exercise}{exercise}{exercises}
\crefname{question}{question}{questions}
\crefname{remark}{remark}{remarks}
\crefname{observation}{observation}{observations}
\crefname{convention}{convention}{conventions}
\crefname{notation}{notation}{notations}

% --------------------------------------------------------------------------
% Bibliography
% --------------------------------------------------------------------------

\usepackage[
  backend=biber,
  style=alphabetic,
  sorting=nyt,
  maxbibnames=99
]{biblatex}

% The present notes use chapter-level prose references rather than citation
% commands.  This guard keeps the project ready for a future bibliography
% without requiring an empty references.bib file in the current release.
\IfFileExists{references.bib}{\addbibresource{references.bib}}{}

% --------------------------------------------------------------------------
% Index
% --------------------------------------------------------------------------

\usepackage{makeidx}
\makeindex

% --------------------------------------------------------------------------
% Drafting tools
% --------------------------------------------------------------------------

% Margin notes can be used during development:
%     \marginnote{Check this proof.}
\usepackage{marginnote}

\newcommand{\todo}[1]{%
  \marginpar{\raggedright\scriptsize\color{WarningRed}\textbf{TODO:} #1}%
}

% Set to true while drafting to display labels in the margins.
\newif\ifshowlabels
\showlabelsfalse

\ifshowlabels
  \usepackage{showkeys}
\fi

% --------------------------------------------------------------------------
% Miscellaneous formatting
% --------------------------------------------------------------------------

\allowdisplaybreaks
\setcounter{tocdepth}{2}
\setcounter{secnumdepth}{3}

\renewcommand{\qedsymbol}{\(\square\)}

% Prevent isolated lines at page boundaries.
\widowpenalty=10000
\clubpenalty=10000
\displaywidowpenalty=10000


% --------------------------------------------------------------------------
% Encoding, fonts, and typography
% --------------------------------------------------------------------------

\usepackage[T1]{fontenc}
\usepackage[utf8]{inputenc} % Harmless on modern LaTeX; useful on older systems.
\usepackage{lmodern}
\usepackage{microtype}

% --------------------------------------------------------------------------
% Page layout and document structure
% --------------------------------------------------------------------------

\usepackage[
  letterpaper,
  inner=1.25in,
  outer=1.00in,
  top=1.00in,
  bottom=1.10in,
  headheight=15pt
]{geometry}

\usepackage{setspace}
\setstretch{1.05}

\usepackage{emptypage}
\usepackage{fancyhdr}

\pagestyle{fancy}
\fancyhf{}
\fancyhead[LE]{\small\thepage}
\fancyhead[RO]{\small\thepage}
\fancyhead[LO]{\small\nouppercase{\rightmark}}
\fancyhead[RE]{\small\nouppercase{\leftmark}}
\renewcommand{\headrulewidth}{0.4pt}
\renewcommand{\footrulewidth}{0pt}

\fancypagestyle{plain}{
  \fancyhf{}
  \fancyfoot[C]{\small\thepage}
  \renewcommand{\headrulewidth}{0pt}
}

\usepackage{titlesec}

\titleformat{\chapter}[display]
  {\normalfont\bfseries}
  {\Large Chapter \thechapter}
  {0.5ex}
  {\Huge}

\titleformat{name=\chapter,numberless}[display]
  {\normalfont\bfseries}
  {}
  {0pt}
  {\Huge}

\titlespacing*{\chapter}{0pt}{-10pt}{30pt}

% --------------------------------------------------------------------------
% Mathematics
% --------------------------------------------------------------------------

\usepackage{amsmath}
\usepackage{amssymb}
\usepackage{amsthm}
\usepackage{mathtools}
\usepackage{mathrsfs}
\usepackage{bm}
\usepackage{dsfont}

\numberwithin{equation}{chapter}

% --------------------------------------------------------------------------
% Lists, tables, and general utilities
% --------------------------------------------------------------------------

\usepackage{enumitem}
\setlist{
  itemsep=2pt,
  topsep=4pt,
  parsep=0pt
}

\usepackage{booktabs}
\usepackage{array}
\usepackage{longtable}
\usepackage{tabularx}
\usepackage{multicol}
\usepackage{ragged2e}

\usepackage{xcolor}
\usepackage{xparse}
\usepackage{etoolbox}
\usepackage{ifthen}
\usepackage{comment}

% --------------------------------------------------------------------------
% Graphics, diagrams, and commutative diagrams
% --------------------------------------------------------------------------

\usepackage{graphicx}
\graphicspath{{figures/}}

\usepackage{tikz}
\usetikzlibrary{
  arrows.meta,
  calc,
  decorations.markings,
  intersections,
  matrix,
  patterns,
  positioning,
  shapes.geometric
}

\usepackage{tikz-cd}

% Basic TikZ styles for fans, lattices, and polytopes.
\tikzset{
  lattice point/.style={
    circle,
    fill,
    inner sep=1.4pt
  },
  ray/.style={
    -{Latex[length=2.2mm]},
    thick
  },
  fan cone/.style={
    fill=black!8,
    draw=black!55
  },
  polytope/.style={
    thick
  },
  toric divisor/.style={
    thick,
    dashed
  }
}

% --------------------------------------------------------------------------
% Color palette
% --------------------------------------------------------------------------

\definecolor{DictionaryBlue}{HTML}{1F4E79}
\definecolor{SurfaceGreen}{HTML}{3F6B4F}
\definecolor{WarningRed}{HTML}{8B2E2E}
\definecolor{ExampleGold}{HTML}{8A641B}
\definecolor{LightGray}{HTML}{F4F4F4}
\definecolor{PastelYellow}{HTML}{FFF8CC}
\definecolor{PausePurple}{HTML}{6F5A8A}
\definecolor{PredictionTeal}{HTML}{287D7A}
\definecolor{CheatsheetSlate}{HTML}{536878}

% --------------------------------------------------------------------------
% Boxed pedagogical environments
% --------------------------------------------------------------------------

\usepackage[most]{tcolorbox}

\tcbset{
  enhanced,
  breakable,
  boxrule=0.7pt,
  arc=1.5mm,
  left=2mm,
  right=2mm,
  top=1.5mm,
  bottom=1.5mm,
  before skip=8pt,
  after skip=8pt
}

\newtcolorbox{chapteroverview}{
  colback=PastelYellow,
  colframe=ExampleGold!70!black,
  title={Chapter overview},
  fonttitle=\bfseries
}

\newtcolorbox{threedictionaries}{
  colback=DictionaryBlue!4,
  colframe=DictionaryBlue,
  title={The three dictionaries},
  fonttitle=\bfseries
}

\newtcolorbox{surfacelaboratory}{
  colback=SurfaceGreen!5,
  colframe=SurfaceGreen,
  title={Surface laboratory},
  fonttitle=\bfseries
}

\newtcolorbox{warningbox}{
  colback=WarningRed!4,
  colframe=WarningRed,
  title={Warning},
  fonttitle=\bfseries
}

\newtcolorbox{guidingquestion}{
  colback=ExampleGold!5,
  colframe=ExampleGold,
  title={Guiding question},
  fonttitle=\bfseries
}

\newtcolorbox{pauseandthink}{
  colback=PausePurple!5,
  colframe=PausePurple,
  title={Pause and think},
  fonttitle=\bfseries
}

\newtcolorbox{prediction}{
  colback=PredictionTeal!5,
  colframe=PredictionTeal,
  title={Prediction},
  fonttitle=\bfseries
}

\newtcolorbox{dictionarytable}{
  colback=DictionaryBlue!3,
  colframe=DictionaryBlue,
  title={Dictionary summary},
  fonttitle=\bfseries
}

\newtcolorbox{agcheatsheet}{
  colback=CheatsheetSlate!5,
  colframe=CheatsheetSlate,
  title={Algebraic geometry cheatsheet},
  fonttitle=\bfseries
}

% --------------------------------------------------------------------------
% Theorem environments
% --------------------------------------------------------------------------

\theoremstyle{plain}
\newtheorem{theorem}{Theorem}[chapter]
\newtheorem{proposition}[theorem]{Proposition}
\newtheorem{lemma}[theorem]{Lemma}
\newtheorem{corollary}[theorem]{Corollary}

\theoremstyle{definition}
\newtheorem{definition}[theorem]{Definition}
\newtheorem{construction}[theorem]{Construction}
\newtheorem{example}[theorem]{Example}
\newtheorem{exercise}[theorem]{Exercise}
\newtheorem{question}[theorem]{Question}

\theoremstyle{remark}
\newtheorem{remark}[theorem]{Remark}
\newtheorem{observation}[theorem]{Observation}
\newtheorem{convention}[theorem]{Convention}
\newtheorem{notation}[theorem]{Notation}

% Unnumbered variants.
\newtheorem*{theorem*}{Theorem}
\newtheorem*{proposition*}{Proposition}
\newtheorem*{definition*}{Definition}
\newtheorem*{example*}{Example}
\newtheorem*{remark*}{Remark}

% --------------------------------------------------------------------------
% Exercises and solutions
% --------------------------------------------------------------------------

\newcounter{exerciseblock}[chapter]
\renewcommand{\theexerciseblock}{\thechapter.\arabic{exerciseblock}}

\NewDocumentEnvironment{exercises}{O{Exercises}}
{
  \section*{#1}
  \addcontentsline{toc}{section}{#1}
  \begin{enumerate}[label=\textbf{\thechapter.\arabic*.},wide=0pt]
}
{
  \end{enumerate}
}

\newif\ifshowsolutions
\showsolutionsfalse
% To compile instructor solutions, change the previous line to:
% \showsolutionstrue

\NewDocumentEnvironment{solution}{}
{
  \ifshowsolutions
    \par\smallskip
    \noindent\textbf{Solution. }\ignorespaces
  \else
    \expandafter\comment
  \fi
}
{
  \ifshowsolutions
    \par\smallskip
  \else
    \expandafter\endcomment
  \fi
}

% --------------------------------------------------------------------------
% Cross-references
% --------------------------------------------------------------------------

\usepackage{hyperref}

\hypersetup{
  colorlinks=true,
  linkcolor=DictionaryBlue,
  citecolor=SurfaceGreen,
  urlcolor=WarningRed,
  pdfauthor={},
  pdftitle={The Three Dictionaries of Toric Geometry},
  pdfsubject={Notes for Renzo's Fall 2026 course at Colorado State University},
  bookmarksopen=true,
  bookmarksnumbered=true
}

\usepackage[nameinlink,noabbrev,capitalise]{cleveref}

\crefname{theorem}{theorem}{theorems}
\crefname{proposition}{proposition}{propositions}
\crefname{lemma}{lemma}{lemmas}
\crefname{corollary}{corollary}{corollaries}
\crefname{definition}{definition}{definitions}
\crefname{construction}{construction}{constructions}
\crefname{example}{example}{examples}
\crefname{exercise}{exercise}{exercises}
\crefname{question}{question}{questions}
\crefname{remark}{remark}{remarks}
\crefname{observation}{observation}{observations}
\crefname{convention}{convention}{conventions}
\crefname{notation}{notation}{notations}

% --------------------------------------------------------------------------
% Bibliography
% --------------------------------------------------------------------------

\usepackage[
  backend=biber,
  style=alphabetic,
  sorting=nyt,
  maxbibnames=99
]{biblatex}

% The present notes use chapter-level prose references rather than citation
% commands.  This guard keeps the project ready for a future bibliography
% without requiring an empty references.bib file in the current release.
\IfFileExists{references.bib}{\addbibresource{references.bib}}{}

% --------------------------------------------------------------------------
% Index
% --------------------------------------------------------------------------

\usepackage{makeidx}
\makeindex

% --------------------------------------------------------------------------
% Drafting tools
% --------------------------------------------------------------------------

% Margin notes can be used during development:
%     \marginnote{Check this proof.}
\usepackage{marginnote}

\newcommand{\todo}[1]{%
  \marginpar{\raggedright\scriptsize\color{WarningRed}\textbf{TODO:} #1}%
}

% Set to true while drafting to display labels in the margins.
\newif\ifshowlabels
\showlabelsfalse

\ifshowlabels
  \usepackage{showkeys}
\fi

% --------------------------------------------------------------------------
% Miscellaneous formatting
% --------------------------------------------------------------------------

\allowdisplaybreaks
\setcounter{tocdepth}{2}
\setcounter{secnumdepth}{3}

\renewcommand{\qedsymbol}{\(\square\)}

% Prevent isolated lines at page boundaries.
\widowpenalty=10000
\clubpenalty=10000
\displaywidowpenalty=10000


% --------------------------------------------------------------------------
% Encoding, fonts, and typography
% --------------------------------------------------------------------------

\usepackage[T1]{fontenc}
\usepackage[utf8]{inputenc} % Harmless on modern LaTeX; useful on older systems.
\usepackage{lmodern}
\usepackage{microtype}

% --------------------------------------------------------------------------
% Page layout and document structure
% --------------------------------------------------------------------------

\usepackage[
  letterpaper,
  inner=1.25in,
  outer=1.00in,
  top=1.00in,
  bottom=1.10in,
  headheight=15pt
]{geometry}

\usepackage{setspace}
\setstretch{1.05}

\usepackage{emptypage}
\usepackage{fancyhdr}

\pagestyle{fancy}
\fancyhf{}
\fancyhead[LE]{\small\thepage}
\fancyhead[RO]{\small\thepage}
\fancyhead[LO]{\small\nouppercase{\rightmark}}
\fancyhead[RE]{\small\nouppercase{\leftmark}}
\renewcommand{\headrulewidth}{0.4pt}
\renewcommand{\footrulewidth}{0pt}

\fancypagestyle{plain}{
  \fancyhf{}
  \fancyfoot[C]{\small\thepage}
  \renewcommand{\headrulewidth}{0pt}
}

\usepackage{titlesec}

\titleformat{\chapter}[display]
  {\normalfont\bfseries}
  {\Large Chapter \thechapter}
  {0.5ex}
  {\Huge}

\titleformat{name=\chapter,numberless}[display]
  {\normalfont\bfseries}
  {}
  {0pt}
  {\Huge}

\titlespacing*{\chapter}{0pt}{-10pt}{30pt}

% --------------------------------------------------------------------------
% Mathematics
% --------------------------------------------------------------------------

\usepackage{amsmath}
\usepackage{amssymb}
\usepackage{amsthm}
\usepackage{mathtools}
\usepackage{mathrsfs}
\usepackage{bm}
\usepackage{dsfont}

\numberwithin{equation}{chapter}

% --------------------------------------------------------------------------
% Lists, tables, and general utilities
% --------------------------------------------------------------------------

\usepackage{enumitem}
\setlist{
  itemsep=2pt,
  topsep=4pt,
  parsep=0pt
}

\usepackage{booktabs}
\usepackage{array}
\usepackage{longtable}
\usepackage{tabularx}
\usepackage{multicol}
\usepackage{ragged2e}

\usepackage{xcolor}
\usepackage{xparse}
\usepackage{etoolbox}
\usepackage{ifthen}
\usepackage{comment}

% --------------------------------------------------------------------------
% Graphics, diagrams, and commutative diagrams
% --------------------------------------------------------------------------

\usepackage{graphicx}
\graphicspath{{figures/}}

\usepackage{tikz}
\usetikzlibrary{
  arrows.meta,
  calc,
  decorations.markings,
  intersections,
  matrix,
  patterns,
  positioning,
  shapes.geometric
}

\usepackage{tikz-cd}

% Basic TikZ styles for fans, lattices, and polytopes.
\tikzset{
  lattice point/.style={
    circle,
    fill,
    inner sep=1.4pt
  },
  ray/.style={
    -{Latex[length=2.2mm]},
    thick
  },
  fan cone/.style={
    fill=black!8,
    draw=black!55
  },
  polytope/.style={
    thick
  },
  toric divisor/.style={
    thick,
    dashed
  }
}

% --------------------------------------------------------------------------
% Color palette
% --------------------------------------------------------------------------

\definecolor{DictionaryBlue}{HTML}{1F4E79}
\definecolor{SurfaceGreen}{HTML}{3F6B4F}
\definecolor{WarningRed}{HTML}{8B2E2E}
\definecolor{ExampleGold}{HTML}{8A641B}
\definecolor{LightGray}{HTML}{F4F4F4}
\definecolor{PastelYellow}{HTML}{FFF8CC}
\definecolor{PausePurple}{HTML}{6F5A8A}
\definecolor{PredictionTeal}{HTML}{287D7A}
\definecolor{CheatsheetSlate}{HTML}{536878}

% --------------------------------------------------------------------------
% Boxed pedagogical environments
% --------------------------------------------------------------------------

\usepackage[most]{tcolorbox}

\tcbset{
  enhanced,
  breakable,
  boxrule=0.7pt,
  arc=1.5mm,
  left=2mm,
  right=2mm,
  top=1.5mm,
  bottom=1.5mm,
  before skip=8pt,
  after skip=8pt
}

\newtcolorbox{chapteroverview}{
  colback=PastelYellow,
  colframe=ExampleGold!70!black,
  title={Chapter overview},
  fonttitle=\bfseries
}

\newtcolorbox{threedictionaries}{
  colback=DictionaryBlue!4,
  colframe=DictionaryBlue,
  title={The three dictionaries},
  fonttitle=\bfseries
}

\newtcolorbox{surfacelaboratory}{
  colback=SurfaceGreen!5,
  colframe=SurfaceGreen,
  title={Surface laboratory},
  fonttitle=\bfseries
}

\newtcolorbox{warningbox}{
  colback=WarningRed!4,
  colframe=WarningRed,
  title={Warning},
  fonttitle=\bfseries
}

\newtcolorbox{guidingquestion}{
  colback=ExampleGold!5,
  colframe=ExampleGold,
  title={Guiding question},
  fonttitle=\bfseries
}

\newtcolorbox{pauseandthink}{
  colback=PausePurple!5,
  colframe=PausePurple,
  title={Pause and think},
  fonttitle=\bfseries
}

\newtcolorbox{prediction}{
  colback=PredictionTeal!5,
  colframe=PredictionTeal,
  title={Prediction},
  fonttitle=\bfseries
}

\newtcolorbox{dictionarytable}{
  colback=DictionaryBlue!3,
  colframe=DictionaryBlue,
  title={Dictionary summary},
  fonttitle=\bfseries
}

\newtcolorbox{agcheatsheet}{
  colback=CheatsheetSlate!5,
  colframe=CheatsheetSlate,
  title={Algebraic geometry cheatsheet},
  fonttitle=\bfseries
}

% --------------------------------------------------------------------------
% Theorem environments
% --------------------------------------------------------------------------

\theoremstyle{plain}
\newtheorem{theorem}{Theorem}[chapter]
\newtheorem{proposition}[theorem]{Proposition}
\newtheorem{lemma}[theorem]{Lemma}
\newtheorem{corollary}[theorem]{Corollary}

\theoremstyle{definition}
\newtheorem{definition}[theorem]{Definition}
\newtheorem{construction}[theorem]{Construction}
\newtheorem{example}[theorem]{Example}
\newtheorem{exercise}[theorem]{Exercise}
\newtheorem{question}[theorem]{Question}

\theoremstyle{remark}
\newtheorem{remark}[theorem]{Remark}
\newtheorem{observation}[theorem]{Observation}
\newtheorem{convention}[theorem]{Convention}
\newtheorem{notation}[theorem]{Notation}

% Unnumbered variants.
\newtheorem*{theorem*}{Theorem}
\newtheorem*{proposition*}{Proposition}
\newtheorem*{definition*}{Definition}
\newtheorem*{example*}{Example}
\newtheorem*{remark*}{Remark}

% --------------------------------------------------------------------------
% Exercises and solutions
% --------------------------------------------------------------------------

\newcounter{exerciseblock}[chapter]
\renewcommand{\theexerciseblock}{\thechapter.\arabic{exerciseblock}}

\NewDocumentEnvironment{exercises}{O{Exercises}}
{
  \section*{#1}
  \addcontentsline{toc}{section}{#1}
  \begin{enumerate}[label=\textbf{\thechapter.\arabic*.},wide=0pt]
}
{
  \end{enumerate}
}

\newif\ifshowsolutions
\showsolutionsfalse
% To compile instructor solutions, change the previous line to:
% \showsolutionstrue

\NewDocumentEnvironment{solution}{}
{
  \ifshowsolutions
    \par\smallskip
    \noindent\textbf{Solution. }\ignorespaces
  \else
    \expandafter\comment
  \fi
}
{
  \ifshowsolutions
    \par\smallskip
  \else
    \expandafter\endcomment
  \fi
}

% --------------------------------------------------------------------------
% Cross-references
% --------------------------------------------------------------------------

\usepackage{hyperref}

\hypersetup{
  colorlinks=true,
  linkcolor=DictionaryBlue,
  citecolor=SurfaceGreen,
  urlcolor=WarningRed,
  pdfauthor={},
  pdftitle={The Three Dictionaries of Toric Geometry},
  pdfsubject={Notes for Renzo's Fall 2026 course at Colorado State University},
  bookmarksopen=true,
  bookmarksnumbered=true
}

\usepackage[nameinlink,noabbrev,capitalise]{cleveref}

\crefname{theorem}{theorem}{theorems}
\crefname{proposition}{proposition}{propositions}
\crefname{lemma}{lemma}{lemmas}
\crefname{corollary}{corollary}{corollaries}
\crefname{definition}{definition}{definitions}
\crefname{construction}{construction}{constructions}
\crefname{example}{example}{examples}
\crefname{exercise}{exercise}{exercises}
\crefname{question}{question}{questions}
\crefname{remark}{remark}{remarks}
\crefname{observation}{observation}{observations}
\crefname{convention}{convention}{conventions}
\crefname{notation}{notation}{notations}

% --------------------------------------------------------------------------
% Bibliography
% --------------------------------------------------------------------------

\usepackage[
  backend=biber,
  style=alphabetic,
  sorting=nyt,
  maxbibnames=99
]{biblatex}

% The present notes use chapter-level prose references rather than citation
% commands.  This guard keeps the project ready for a future bibliography
% without requiring an empty references.bib file in the current release.
\IfFileExists{references.bib}{\addbibresource{references.bib}}{}

% --------------------------------------------------------------------------
% Index
% --------------------------------------------------------------------------

\usepackage{makeidx}
\makeindex

% --------------------------------------------------------------------------
% Drafting tools
% --------------------------------------------------------------------------

% Margin notes can be used during development:
%     \marginnote{Check this proof.}
\usepackage{marginnote}

\newcommand{\todo}[1]{%
  \marginpar{\raggedright\scriptsize\color{WarningRed}\textbf{TODO:} #1}%
}

% Set to true while drafting to display labels in the margins.
\newif\ifshowlabels
\showlabelsfalse

\ifshowlabels
  \usepackage{showkeys}
\fi

% --------------------------------------------------------------------------
% Miscellaneous formatting
% --------------------------------------------------------------------------

\allowdisplaybreaks
\setcounter{tocdepth}{2}
\setcounter{secnumdepth}{3}

\renewcommand{\qedsymbol}{\(\square\)}

% Prevent isolated lines at page boundaries.
\widowpenalty=10000
\clubpenalty=10000
\displaywidowpenalty=10000
